\section{Construct a BST from Preorder Traversal}
\label{sec:bst-construct-preorder}

\problemstatement{Given an array representing the preorder traversal of a
BST (root, then left subtree, then right subtree, with no duplicate
values), reconstruct the BST and return its root.}

\begin{patternbox}
\emph{``Reconstruct a BST from one traversal''} works only because the
traversal is specifically \textbf{preorder} \emph{and} the tree is
specifically a \textbf{BST} -- the very first element is always the root,
and the BST invariant (combined with knowing each value's relationship to
the root) is enough to tell us, for every subsequent element, whether it
belongs in the left subtree or the right, without needing a second
traversal (like inorder) to disambiguate, the way general binary tree
reconstruction usually requires.
\end{patternbox}

\subsection*{Understanding the problem carefully}

In a preorder traversal of a BST, the first value is the root. Every
value immediately following it that is \textbf{less than} the root
belongs to the root's \emph{left} subtree -- and because of how preorder
works (left subtree fully traversed before any of the right subtree
begins), \emph{all} such smaller values appear contiguously, immediately
after the root, before any larger value appears. The first value that is
\textbf{greater} than the root marks the start of the right subtree, and
every later value will also be greater than the root.

\begin{ideabox}[Build with an Upper Bound, Consuming a Shared Pointer]
Maintain a single index pointer into the preorder array, shared across
all recursive calls (advanced every time a node is consumed). Define
$\textsc{Build}(\textit{upperBound})$: if the pointer is exhausted, or
the next value to consume exceeds $\textit{upperBound}$, return null
(this value belongs to some ancestor's other subtree, not here). Otherwise,
consume the next value as a new node; recursively build its left subtree
with $\textit{upperBound}$ tightened to this new node's value; then
recursively build its right subtree with the \emph{original}
$\textit{upperBound}$ unchanged (everything remaining that is still
$\le$ the original bound, but now also $>$ this node's value, belongs on
the right).
\end{ideabox}

\subsection*{Why a single shared pointer and a bound (not an explicit split search) suffices}

The key insight is that we never need to explicitly search for ``where
the left subtree ends and the right subtree begins'' -- the recursive
calls naturally stop consuming values the moment they encounter one that
violates the current $\textit{upperBound}$, which is exactly the boundary
between subtrees. Because the shared pointer only ever advances (never
resets or searches backward), each value in the preorder array is
considered, and consumed, exactly once across the entire construction,
regardless of how deeply nested its eventual position in the tree is.

\subsection*{Algorithm}

\begin{enumerate}
  \item Initialize a shared index $i \gets 0$.
  \item $\textsc{Build}(\textit{upperBound})$:
  \begin{enumerate}
    \item If $i \ge n$ or $\textit{preorder}[i] > \textit{upperBound}$:
          return null.
    \item Create a new node with value $\textit{preorder}[i]$;
          increment $i$.
    \item Recursively set this node's left child to
          $\textsc{Build}(\textit{node.val})$.
    \item Recursively set this node's right child to
          $\textsc{Build}(\textit{upperBound})$ (the original bound).
    \item Return this node.
  \end{enumerate}
  \item Call $\textsc{Build}(+\infty)$ at the top level.
\end{enumerate}

\subsection*{Dry run}

\dryrun{Take $\textit{preorder} = [8,5,1,7,10,12]$.}

\begin{itemize}
  \item $\textsc{Build}(+\infty)$: consume $8$ ($i \to 1$). Node $8$.
  \item Left: $\textsc{Build}(8)$: consume $5$ ($5\le8$, $i\to2$). Node
        $5$.
  \begin{itemize}
    \item Left: $\textsc{Build}(5)$: consume $1$ ($1\le5$, $i\to3$).
          Node $1$. Left: $\textsc{Build}(1)$: next is $7>1$, return
          null. Right: $\textsc{Build}(5)$: next is $7>5$, return null.
          Node $1$ is a leaf.
    \item Right: $\textsc{Build}(8)$: consume $7$ ($7\le8$, $i\to4$).
          Node $7$. Left: $\textsc{Build}(7)$: next is $10>7$, return
          null. Right: $\textsc{Build}(8)$: next is $10>8$, return null.
          Node $7$ is a leaf.
  \end{itemize}
  \item Right of $8$: $\textsc{Build}(+\infty)$: consume $10$
        ($i\to5$). Node $10$.
  \begin{itemize}
    \item Left: $\textsc{Build}(10)$: next is $12>10$, return null.
    \item Right: $\textsc{Build}(+\infty)$: consume $12$ ($i\to6$).
          Node $12$, a leaf.
  \end{itemize}
\end{itemize}

Resulting tree: root $8$, left child $5$ (children $1,7$), right child
$10$ (right child $12$).

\subsection*{C++ implementation}

\begin{lstlisting}
#include <bits/stdc++.h>
using namespace std;

struct TreeNode {
    int val;
    TreeNode* left;
    TreeNode* right;
    TreeNode(int v) : val(v), left(nullptr), right(nullptr) {}
};

TreeNode* build(vector<int>& preorder, int& i, int upperBound) {
    if (i >= (int)preorder.size() || preorder[i] > upperBound) {
        return nullptr;
    }
    TreeNode* node = new TreeNode(preorder[i++]);
    node->left = build(preorder, i, node->val);
    node->right = build(preorder, i, upperBound);
    return node;
}

TreeNode* bstFromPreorder(vector<int>& preorder) {
    int i = 0;
    return build(preorder, i, INT_MAX);
}

int main() {
    vector<int> preorder = {8, 5, 1, 7, 10, 12};
    TreeNode* root = bstFromPreorder(preorder);
    cout << "Root: " << root->val << endl;
    return 0;
}
\end{lstlisting}

\begin{complexitybox}
\textbf{Time Complexity.} Each element of the preorder array is consumed
exactly once, with $O(1)$ work per element, giving
\[
  O(n).
\]
\textbf{Space Complexity.} $O(h)$ for the recursion stack, plus $O(n)$
for the $n$ nodes created.
\end{complexitybox}

\begin{takeawaybox}
Reconstructing a tree from a single traversal is only possible when that
traversal, combined with some known structural invariant (here, the BST
property), removes all ambiguity about subtree boundaries. The ``pass a
bound downward, consume from a shared pointer'' technique used here is a
direct cousin of Validate BST's (Section~\ref{sec:bst-validate}) bound-passing
recursion, repurposed for construction instead of verification.
\end{takeawaybox}
